\section{A uniform finite engine for rational complex polynomials}
\label{sec:polynomial-algebra}

The minimum-modulus argument requires factorization of a rational complex
polynomial whose degree varies with the requested accuracy. The repository's
FTA draft only asserts its full-factorization schema for each fixed degree.
That schema does not supply the uniform construction needed here. We give a
separate finite construction for rational polynomial codes. Its only
topological ingredient is the elementary winding calculus for rational
polygons: additivity under subdivision, invariance under a homotopy with a
positive separation margin, and winding one for a positively oriented square
about an interior point. These finite facts do not use the fundamental
theorem of algebra, Jensen's formula, or a general argument principle.

\subsection{Coefficient bounds and the squarefree reduction}

For a rational or Gaussian-rational polynomial $A(X)=\sum a_jX^j$, put
\[
 \mathcal B(A,B)=1+\sum_j |a_j|_1B^j,\qquad
 \mathcal B_r(A,B)=1+\sum_{j\ge r}
       j(j-1)\cdots(j-r+1)|a_j|_1B^{j-r},
\]
where $B\ge1$ is rational and $|x+iy|_1=|x|+|y|$.
These are rational upper bounds for the polynomial and its indicated
derivatives on $|z|\le B$. All ensuing constants can be computed exactly;
there is no appeal to an unspecified bound on coefficient heights.

The zero polynomial is rejected by the input guard; a nonzero constant
returns an empty root list and itself as leading coefficient. For a
nonconstant input first divide by its nonzero leading coefficient $a$,
so that the polynomial used in the gcd loop is monic.
Euclidean division and the extended Euclidean algorithm over $\mathbb Q(i)$
are bounded algorithms: a division has at most the current degree plus one
steps, and a nonzero Euclidean remainder strictly decreases the degree.
Leading coefficients are rational complex values, so their nonvanishing is
decidable and their inverses are exact rational operations. The usual gcd
identities give the squarefree decomposition
\[
                  P=a\prod_{j=1}^{d}S_j^{\,j},
\]
with monic, squarefree, pairwise coprime $S_j$, allowing $S_j=1$.
For clarity, writing $P_0=P/a$, this can be computed by setting
$c=\gcd(P_0,P_0')$, $w=P_0/c$ and,
successively, $y=\gcd(w,c)$, $S_j=w/y$, $w=y$, $c=c/y$.
The loop stops when $w=1$ and fills remaining entries with $1$; a zero
Euclidean remainder stops that particular gcd computation. The outer
loop runs at most $d=\deg P$ times. Polynomial identities and the
degree decrease prove the decomposition; its implementation can verify the
displayed product directly. Thus it suffices to treat a monic squarefree
polynomial $S$ of degree $d\ge1$.

Compute polynomials $U,V$ with
\begin{equation}
                       US+VS'=1.                    \label{eq:poly-bezout}
\end{equation}
Let $R$ be an integer greater than $2+\sum_{j<d}|s_j|_1$, and set
$B=4(R+1)$. Leading-term dominance gives $|S(z)|\ge1$ whenever
$|z|\ge R-1$. For a root name $\alpha$, read a rational approximation
$a$ with error at most $1/4$ and decide the rational comparison
$|a|\le R-1/2$ by squaring. The opposite branch would imply
$|\alpha|>R-3/4>R-1$ and contradict the displayed lower bound.
Consequently $|\alpha|\le R-1/4<R$. This obtains a strict bound with
a margin, rather than inferring it from the negation of an order relation.
Put
\[
 A=\mathcal B(U,B),\quad K=\mathcal B(V,B),\quad
 M_1=\mathcal B_1(S,B),\quad M_2=\mathcal B_2(S,B).
\]
In particular $A,K,M_1,M_2\ge1$. At a root $\alpha$,
\begin{equation}
                  |S'(\alpha)|\ge K^{-1}.            \label{eq:poly-root-derivative}
\end{equation}
The finite Taylor identity gives
\[
 |S(\beta)-S(\alpha)-S'(\alpha)(\beta-\alpha)|
       \le\tfrac12 M_2|\beta-\alpha|^2.
\]
Consequently two roots at distance at most $(KM_2)^{-1}$ must be equal:
the last inequality and \eqref{eq:poly-root-derivative} imply
$|\beta-\alpha|/K\le |\beta-\alpha|/(2K)$.
This proves the needed separation for roots known to be apart without
dividing by an undecided difference of named numbers.

\subsection{Real roots by a bounded grid and bisection}

The following real procedure will be used several times. Given a nonzero
rational polynomial $F$, remove repeated factors by gcd, normalize to monic,
and call the result $F_0$. Compute $U_0F_0+V_0F_0'=1$, a Cauchy root bound
$R_0$, and bounds $K_0\ge |V_0|$, $M_0\ge\max(1,|F_0''|)$ on
$[-R_0-1,R_0+1]$. Choose
\[
                       h=\frac1{64K_0M_0}.
\]
For grid centres $jh$ covering $[-R_0-h,R_0+h]$, inspect the overlapping
intervals $[jh-2h,jh+2h]$. Retain an interval precisely when its two rational
endpoint values of $F_0$ have strictly opposite signs. This is a finite
scan with fewer than $2R_0/h+10$ centres.

Every retained interval contains a uniquely specified real root name.
Bisection uses exact rational signs; an exactly zero midpoint produces a
constant name. After $q$ bisections the width is $4h2^{-q}$, so the number
of bisections for any requested accuracy is explicitly bounded. The nested
midpoints are Cauchy, and the polynomial Lipschitz bound shows their
residual tends to zero at the same explicit rate. Existence here is this
construction, rather than an invocation of a general intermediate-value
principle.

Conversely, every real root $a$ of $F_0$ lies in one of the retained
intervals. Choose a grid centre within $h$ of an approximation to $a$,
with the approximation fine enough to keep the distance below $3h/2$.
On that interval the derivative differs from $F_0'(a)$ by less than
$1/(4K_0)$, while $|F_0'(a)|\ge1/K_0$. The finite Taylor identity therefore
makes the endpoint signs strictly opposite. No exact comparison of $a$
with a mesh point is required. Distinct roots are at least
$(K_0M_0)^{-1}$ apart; any two overlapping retained intervals consequently
name the same root. Merge overlapping intervals and keep one representative
per group. This gives a finite list of pairwise apart root codes containing
every real root. Each code consists of $F_0$ and one strict sign-changing
rational interval. Its evaluator is the fixed bisection program just
described, uniform in the numerical polynomial code and the precision.

The identities $F_0\mid F$ and $F\mid F_0^{\deg F}$, obtainable during
the gcd calculation, prove that the list also describes precisely the
real zeros of the original $F$.

\subsection{Rational coordinate polynomials for complex roots}

Set $Q=S\overline S\in\mathbb Q[X]$, where conjugation acts on the
coefficients, and let $n=2d$. Consider the finite rational algebra represented
by arrays in the basis $u^av^b$, $0\le a,b<n$, with multiplication reduced
by the monic relations $Q(u)=Q(v)=0$. Its array dimension is $n^2$.
Form the multiplication matrices for
\[
                        x_0=(u+v)/2,
             \qquad     y_0^2=-(u-v)^2/4.
\]
Let $A_X$ be the characteristic polynomial of the first matrix and let
$A_Y(t)$ be the characteristic polynomial of the second matrix evaluated
at $t^2$. Both are nonzero monic rational polynomials, of degrees $n^2$
and $2n^2$, respectively. Determinants can be computed by the finite
permutation formula; efficiency is unnecessary here.

If $\alpha=x+iy$ is a root of $S$, then $Q(\alpha)=Q(\overline\alpha)=0$.
Evaluating the array algebra at $(u,v)=(\alpha,\overline\alpha)$ and using
Cayley--Hamilton gives
\begin{equation}
                         A_X(x)=A_Y(y)=0.              \label{eq:coordinate-polynomials}
\end{equation}
The use of Cayley--Hamilton is purely finite: the adjugate identity for
$tI-M$ gives its characteristic-polynomial identity, and reduction by the
two monic polynomials gives an identity of coefficient arrays. No
existence of complex roots has been used to construct these matrices.

Apply the real-root procedure to $A_X,A_Y$. Every possible complex root
of $S$ is now among a finite Cartesian product of coded real algebraic
numbers. The remaining task is an exact, bounded zero test for each pair.

\subsection{Exact zero and order tests on these codes}

Suppose $x,y$ are real root codes for monic rational polynomials $F,G$,
and $W(X,Y)$ is a rational polynomial. In the finite array algebra
\[
                  \mathbb Q[X,Y]/(F(X),G(Y)),
\]
compute the characteristic polynomial $C$ of multiplication by $W$.
It is monic and satisfies $C(W(x,y))=0$. Factor out its initial power of
the variable by examining its rational coefficients:
\[
                    C(t)=t^mC_0(t),\qquad C_0(0)\ne0.
\]
If $C_0$ is constant, $W(x,y)=0$. Otherwise write $C_0(t)=\sum c_jt^j$
and choose the positive rational
\begin{equation}
 \eta=\min\left(1,\frac{|c_0|}{2(1+\sum_{j\ge1}|c_j|)}\right).
                                                        \label{eq:algebraic-gap}
\end{equation}
Every nonzero real root of $C$ has absolute value at least $\eta$.
Indeed, for $|t|<\eta$, the finite geometric estimate gives
$|C_0(t)|\ge |c_0|/2$; its reciprocal is then permitted, and
$C(t)=0$ implies $t^m=0$ and hence $t=0$.

Evaluate $W(x,y)$ with error at most $\eta/8$, using the bisection codes
and the coefficient bound for the two-variable polynomial on their known
input boxes. This fixes a numerical precision before any evaluation loop.
For the rational result $w$, the test $|w|\le\eta/2$ means the exact
value is zero; otherwise the value is apart from zero and has the sign of
$w$. The gap \eqref{eq:algebraic-gap} proves both branches. The zero
branch has a proof at every precision by the displayed annihilating
identity, rather than by a real-equality oracle.

Apply the test to the real and imaginary parts of $S(x+iy)$. Retain
exactly the pairs passing both zero tests. Equations
\eqref{eq:coordinate-polynomials} and completeness of the real-root lists
prove that this is a complete list of actual roots. Different retained
pairs are apart because their real-root intervals distinguish at least
one coordinate. The same test, applied to $x^2+y^2-r^2$, decides whether
a root is inside, on, or outside a rational circle. Thus a later selection
of roots with $|\alpha|\le3$ includes boundary roots correctly.

\subsection{Why the complete list has exactly the degree}

The preceding construction lists all roots but has not assumed that the
list is nonempty. The following quantitative argument proves its length
is $d$, and closes that possible circularity.

We first need a bounded local root construction. Let $z$ be Gaussian
rational with $|z|\le B-1$, and suppose $|S(z)|\le\varepsilon$, where
\begin{equation}
 \varepsilon\le\min\left(
       \frac1{4A},\frac1{8K},\frac1{16K^2M_2},
       \frac b{16K}\right),\qquad 0<b\le1.              \label{eq:newton-budget}
\end{equation}
By \eqref{eq:poly-bezout}, $|S'(z)|\ge1/(2K)$. Set
\[
        r=4K\varepsilon,\qquad
        T(w)=w-S(w)/S'(z).
\]
The denominator is a nonzero Gaussian rational and is computed exactly.
On $|w-z|\le r$, the finite difference-of-powers formula gives
\[
 |T(w)-T(w')|\le2KM_2r\,|w-w'|\le\tfrac12|w-w'|,
       \qquad |T(z)-z|\le r/2.
\]
Thus the rational iteration $z,T(z),T^2(z),\ldots$ remains in that
disk and has a geometric Cauchy modulus. It gives a root name within
$r\le b/4$ of $z$. Its step count is the requested precision plus a
computed binary length; this is one primitive recursion, not an unbounded
Newton search. The coefficient identity also gives a geometric residual
bound for $S$ at the iterates.

Now put $r_0=(64KM_2)^{-1}$. Around every root in the complete list choose
a rational square of half-side $r_0$, whose centre approximates that root
with error at most $r_0/16$. These squares are disjoint, by root separation.
Their boundaries have distance at least $15r_0/16$ from their corresponding
roots. The Taylor inequality and \eqref{eq:poly-root-derivative} give a
zero-free homotopy of $S$ on each boundary to
$S'(\alpha)(z-\alpha)$: the remainder is less than $1/32$ of the linear
term. The image therefore has winding one.

Take $b=r_0/2$ in \eqref{eq:newton-budget}. At every rational point outside
the interiors of all these squares and inside the outer square of
half-side $R+1$, one has
\begin{equation}
                              |S(z)|>\varepsilon.       \label{eq:poly-outside-lower}
\end{equation}
Otherwise the just-proved rational iteration would construct a root within
$b/4$ of $z$. Completeness of the algebraic list identifies it with a listed
root, whose distance from every such $z$ is greater than $b$, a contradiction.
This conclusion uses a decidable rational comparison: the branch
$|S(z)|\le\varepsilon$ is eliminated by the explicit construction.

Partition the outer square minus the root squares by a rational rectangular
grid, including every square edge and further subdividing until every cell
has diameter at most $\varepsilon/(4M_1)$. All cutoffs are explicit. The
images of each cell's vertices, and the polynomial images of its edges,
stay in a disk of radius at most $\varepsilon/4$ about the image of one
vertex. By \eqref{eq:poly-outside-lower} that disk misses zero. Each cell
image has winding zero. Cancelling shared edges leaves
\[
 \operatorname{wind}(S\circ\partial Q_{R+1})
       =\sum_{\text{root squares}}1.
\]
On the outer square, the leading term of the monic $S$ dominates the sum
of the lower terms. The zero-free straight homotopy to $z^d$ gives winding
$d$. This proves that the list contains exactly $d$ roots.

Synthetic division now yields
\[
                             S(z)=\prod_{j=1}^d(z-\alpha_j).
\]
When a root is removed, every remaining root is still a root of the
quotient, since its difference from the removed root has an explicit
positive lower bound. Repeating the argument for the rational squarefree
factors $S_j$ gives the full multiset factorization of the original $P$,
including multiplicities.

\subsection{Primitive recursion and the induction audit}

Every output root is evaluated by the same fixed bounded bisection routine
on its finite polynomial-and-interval code. Full factorization therefore
does not iterate a higher-order root-selection operation of variable depth.
The operations are rational arithmetic, coefficient-array reduction,
degree-bounded Euclidean division, determinant expansion over a finite
permutation list, bounded grid scans, bisection with a specified number of
steps, and a geometrically bounded rational iteration. Their loop lengths
are obtained by composition, primitive recursion, factorial, powers,
rational ceilings, and bounded binary logarithms. Computed coefficient
bounds replace any hidden root-separation or height oracle.

For the internal proof, requested precisions are fixed before induction.
The invariants are the exact division identities, degree decreases,
matrix-array identities, rational brackets and polynomial signs, the
rational contraction and residual inequalities, and the finite signed-edge
balance. Bounded lists of such conditions are encoded by Boolean folds.
Synthetic-division identities and products of algebraic root names are
checked at an allocated finite precision, using root bounds and the finite
telescoping product inequality; all root evaluations use their uniform
finite codes. These are quantifier-free induction arguments.

This establishes the uniform finite polynomial engine required below.
It does not postulate a total evaluator for arbitrary primitive recursive
programs, a compactness principle, a general real zero test, or the
fundamental theorem of algebra as an unproved analytic premise.
